\documentclass{letask}
\usepackage{textcomp}	%типографские знаки

\begin{document}
\include{cover}	
	
\begin{flushright}
	\textit{}\\
	\textit{}\\
\end{flushright}

% \begin{center}
% 	\Large \bfseries {Лабораторная работа \textnumero{4.1.3}}\\	
% 	\huge \bfseries РЕФРАКТОМЕТР АББЕ
% \end{center}

\textbf{Цель работы:} Знакомство с методом измерения показателей прелом- ления твёрдых и жидких сред в монохроматическом свете.

\textbf{В работе используются:} технический рефрактометр Аббе; освети- тель; набор стеклянных образцов; жидкости с неизвестными показате- лями преломления (глицерин, этиловый спирт); монобромнафталин; дистиллированная вода.

\section*{Теоретические сведения}

Технический рефрактометр Аббе служит для быстрого (и сравни- тельно грубого) измерения показателей преломления жидких и твёрдых тел. Устройство рефрактометра Аббе основано на явлении полного внут- реннего отражения.

% \begin{figure}[h!]
% 	\begin{center}
% 		\includegraphics[width =  0.6\lw]{mods}
% 		\caption{Предельный угол полного внутреннего отражения  (а) и предельный угол преломления (б)}
% 		\label{spectr}
% 	\end{center}	
% \end{figure}


\section*{Экспериментальная установка}

% \begin{figure}[H]
% \centering
% 	\begin{center}
% 		\includegraphics[width = 0.9 \lw]{img2}
% 	\end{center}
% 	\caption{Схема экспериментальной установки}
% \end{figure}



\newpage
\section*{Ход работы}


\section*{Вывод}

\begin{itemize}
    \item При измерении зависимости видности от угла поляризации (при нулевой разности хода) единичная видность не была достигнута (при $0^{\circ}$ она составила $\nu_{3\text{, max}} = 0.72$). Это можно объяснить неидеальной когерентностью лазерных лучей из-за их дифракции на дефектах оптических элементов (царапины, пылинки), а также неточной установкой зеркала $\text{З}_2$ на положение нулевой разности хода. 
    \item Ближе всего к линейной зависимости -- зависимость от первой степени косинуса угла поляризации, из чего можно сделать вывод, что поляризация лазерного излучения скорее линейная с изменяющимся в небольших пределах направлением поляризации.
    \item При измерении зависимости видности от разности хода лучей (при одинаковой поляризации лучей) были оценены: размер резонатора, межмодовое расстояние, диапазон генерации частот и число генерируемых мод.
\end{itemize}

\end{document}
